\documentclass[11pt]{article}
\usepackage{amsmath,amssymb,amsthm}
\usepackage{hyperref}

\newtheorem{theorem}{Theorem}
\newtheorem{proposition}[theorem]{Proposition}
\newtheorem{lemma}[theorem]{Lemma}
\newtheorem{corollary}[theorem]{Corollary}
\newtheorem{remark}[theorem]{Remark}

\newcommand{\Z}{\mathbb{Z}}
\newcommand{\qpochn}[3]{(#1;#2)_{#3}}
\newcommand{\qbin}[2]{\genfrac{[}{]}{0pt}{}{#1}{#2}}

\title{Corrected $q$-Binomial Inverse for $h_m$ and $Q_n$}
\author{Notes (February 2026)}
\date{}

\begin{document}
\maketitle

\section{The corrected inverse formula}

The forward $q$-binomial relation between $Q_n$ and $h_m$ is:
\begin{equation}\label{eq:forward}
  Q_n = \sum_{j=0}^{n} (-1)^j\, q^{\binom{j+1}{2}}\,
    \qbin{n}{j}\, h_{n-j}.
\end{equation}

\begin{proposition}\label{prop:inverse}
The inverse of \eqref{eq:forward} is:
\begin{equation}\label{eq:inverse}
  h_m = \sum_{k=0}^{m} q^k\, \qbin{m}{k}\, Q_{m-k}.
\end{equation}
\end{proposition}

\begin{proof}
From the E--H factorisation $\mathrm{GK}_c(z,q) = H(z) \cdot F_c(z,q)$
where $H(z) = 1/\qpochn{zq}{q}{\infty}$ and
$F_c(z,q) = \qpochn{zq}{q}{\infty} \cdot \mathrm{GK}_c(z,q)$, we have:
\begin{align*}
  [z^m]\, H(z) &= \frac{q^m}{\qpochn{q}{q}{m}}, \\
  [z^m]\, F_c(z,q) &= \frac{Q_m}{\qpochn{q}{q}{m}}.
\end{align*}
The first identity is Euler's partition identity
$1/\qpochn{z}{q}{\infty} = \sum_{m \geq 0} z^m/\qpochn{q}{q}{m}$
evaluated at $z \mapsto zq$, giving $[z^m] = q^m/\qpochn{q}{q}{m}$.

Since $\mathrm{GK}_c = H \cdot F_c$, extracting $[z^m]$ gives:
\[
  g_m = \sum_{k=0}^{m} [z^k]\,H(z) \cdot [z^{m-k}]\,F_c(z,q)
  = \sum_{k=0}^{m} \frac{q^k}{\qpochn{q}{q}{k}} \cdot
    \frac{Q_{m-k}}{\qpochn{q}{q}{m-k}}.
\]
Multiplying both sides by $\qpochn{q}{q}{m}$ and using
$h_m = \qpochn{q}{q}{m} \cdot g_m$:
\[
  h_m = \sum_{k=0}^{m} \frac{\qpochn{q}{q}{m}}
    {\qpochn{q}{q}{k}\,\qpochn{q}{q}{m-k}}
    \cdot q^k \cdot Q_{m-k}
  = \sum_{k=0}^{m} q^k\, \qbin{m}{k}\, Q_{m-k}. \qedhere
\]
\end{proof}

\begin{remark}
An earlier version of \texttt{definitions.tex} had $q^{j^2}$ in place
of $q^j$ in the inverse formula.  The error likely arose from confusion
with the forward formula's $q^{\binom{j+1}{2}}$.  The key difference:
\begin{itemize}
  \item Forward: the sign-reversing Euler factor
    $\qpochn{zq}{q}{\infty} = \sum (-1)^j q^{\binom{j+1}{2}} z^j / \qpochn{q}{q}{j}$
    contributes $q^{\binom{j+1}{2}}$.
  \item Inverse: the partition-counting factor
    $1/\qpochn{zq}{q}{\infty} = \sum q^j z^j / \qpochn{q}{q}{j}$
    contributes $q^j$.
\end{itemize}
\end{remark}

\section{Immediate corollary: $Q \geq 0$ implies $h \geq 0$}

\begin{corollary}
If $Q_n \in \Z_{\geq 0}[q]$ for all $n \geq 0$, then
$h_m \in \Z_{\geq 0}[q]$ for all $m \geq 0$.
\end{corollary}

\begin{proof}
Every term in \eqref{eq:inverse} is a product of non-negative
polynomials: $q^k \geq 0$, $\qbin{m}{k} \in \Z_{\geq 0}[q]$,
and $Q_{m-k} \in \Z_{\geq 0}[q]$ by hypothesis.
\end{proof}

The converse fails: $h_m \geq 0$ does not imply $Q_n \geq 0$,
because the forward formula \eqref{eq:forward} is an alternating sum.

\section{Corrected examples for $d = 2$, $c = (1,1,0)$}

With $Q_n = q^{n^2}$ (first Rogers--Ramanujan), the inverse gives:
\[
  h_m = \sum_{k=0}^{m} q^{k + (m-k)^2}\, \qbin{m}{k}.
\]
\begin{align*}
  h_0 &= 1, \\
  h_1 &= q^{0+1}\qbin{1}{0} + q^{1+0}\qbin{1}{1} = q + q = 2q, \\
  h_2 &= q^{0+4}\qbin{2}{0} + q^{1+1}\qbin{2}{1} + q^{2+0}\qbin{2}{2}
       = q^4 + q^2(1+q) + q^2 = 2q^2 + q^3 + q^4.
\end{align*}

\section{Dominance investigation (negative result)}

We explored whether $Q_n \geq 0$ could be explained by
coefficient-wise dominance relations between successive $h_m$.
Define:
\begin{align*}
  D_m &= h_m - q \cdot h_{m-1}, \\
  R_n &= \sum_{j=1}^{n} (-1)^{j-1}\, q^{\binom{j+1}{2}}\,
    \qbin{n}{j}\, h_{n-j},
\end{align*}
so that $Q_n = h_n - R_n$.  Computations for
$d = 1, 2, 4, 5$ confirm:
\begin{enumerate}
  \item $D_m \geq 0$ for all $m$ and all tested profiles.
  \item $R_n \geq 0$ for all $n$ and all tested profiles.
  \item Even partial sums $P_n^{(2k)} = \sum_{j=0}^{2k}(\cdots) \geq 0$.
\end{enumerate}

However, all three properties are \textbf{consequences} of $Q_n \geq 0$,
not independent facts:

\begin{proposition}
$R_n = \sum_{k=1}^{n} q^k\, \qbin{n}{k}\, Q_{n-k}$.
\end{proposition}

\begin{proof}
This is the inverse formula \eqref{eq:inverse} with the $k=0$ term
($Q_n$) removed: $R_n = h_n - Q_n = \sum_{k=1}^n q^k \qbin{n}{k} Q_{n-k}$.
\end{proof}

\begin{proposition}
$D_m = Q_m + \sum_{j=1}^{m} q^{2j}\, \qbin{m-1}{j}\, Q_{m-j}$.
\end{proposition}

\begin{proof}
From \eqref{eq:inverse}:
$D_m = \sum_{k=0}^m q^k \qbin{m}{k} Q_{m-k}
     - \sum_{k=0}^{m-1} q^{k+1} \qbin{m-1}{k} Q_{m-1-k}$.
Reindex the second sum ($k \mapsto j-1$) and use
$\qbin{m}{j} - \qbin{m-1}{j-1} = q^j \qbin{m-1}{j}$.
\end{proof}

Thus $D_m \geq 0$ and $R_n \geq 0$ follow from $Q_k \geq 0$ for
$k \leq m$ (resp.\ $k \leq n-1$).
The positivity of $Q_n$ cannot be decomposed into simpler
pairwise dominance inequalities via this approach.

\bigskip
\noindent\textbf{Conclusion.}
The conjecture $Q_n \geq 0$ appears to be irreducibly about the
interaction between the full sequence $\{h_m\}$ (encoding the CPP
structure via $h_m = \qpochn{q}{q}{m} \cdot g_m$) and the
$q$-binomial transform.  Proving it likely requires a direct
construction (bijective or algebraic) rather than reduction to
simpler inequalities.

\end{document}
